\chapter{Εισαγωγή}
\pagenumbering{arabic}%\pagestyle{mythesis}
\label{ch:intro}
\hrule\vspace{0.2in}
\begin{itemize}
  \item[\ref{sec:11}] Στόχοι
  \item[\ref{sec:12}] Δομή της Διατριβής
\end{itemize}
\hrule \vspace{0.2in}



\section{Στόχοι}\label{sec:11}
Εδώ αρχίζει η κανονική αρίθμηση των σελίδων. Παραπάνω η αρίθμηση
είναι i, ii, iii, όπως φαίνεται.

Συνήθως αυτό το κεφάλαιο (ΕΙΣΑΓΩΓΗ) πρέπει να υπάρχει. Στο
περιεχόμενό του πρέπει να συμπεριλαμβάνονται και τα ακόλουθα:

\y Στόχοι. Για παράδειγμα: Ο συνολικός στόχος της διατριβής είναι
Οι επί μέρους στόχοι:
\begin{itemize}
\item Να αναλυθεί η συμπεριφορά ...
\item Η ανάπτυξη  ...
\item Να κατηγοριοποιηθούν  ...
\item Να επιβεβαιωθούν  ...
\end{itemize}




\section{Δομή της Διατριβής}\label{sec:12}
Η διατριβή περιέχει ν κεφάλαια: Το Κεφάλαιο 1.


\newpage
Το Σχήμα~\ref{fig:TwoBirds} δείχνει το Πρώτο Σχήμα του Πρώτου Κεφαλαίου.

\begin{figure}[t]
	\centering
	\begin{subfigure}[t]{0.49\textwidth}
		\centering
		\includegraphics[width=0.25\textwidth]{Figures/bird.eps}
		\caption{}
		\label{subfig:TwoBirdsA}
	\end{subfigure}
	\hfill
	\begin{subfigure}[t]{0.49\textwidth}
		\centering
		\includegraphics[width=0.25\textwidth]{Figures/bird.eps}
		\caption{}
		\label{subfig:TwoBirdsB}
	\end{subfigure}
	\caption{Το σχήμα (a) είναι δίπλα στο σχήμα (b).}
	\label{fig:TwoBirds}
\end{figure}